\chapter*{Préface}
\addcontentsline{toc}{chapter}{Préface}

Ce mémoire a été réalisé dans le cadre du \textbf{[intitulé du master]} de \textbf{[université]}, de \textbf{[mois année]} à \textbf{[mois année]}, au sein de l'IHU SEPSIS, dans le Work Package~6 consacré au sepsis. Il a été encadré par \textbf{[Prénom Nom]}, et sa partie médicale a été validée par le Pr Nicholas Heming, de l'Hôpital Raymond-Poincaré (AP-HP).

Le travail porte sur une question à la fois technique et clinique : peut-on extraire de façon fiable les informations contenues dans des dossiers médicaux scannés de patients septiques, et savoir quand le modèle qui les extrait se trompe ? Les grands modèles de langage multimodaux savent lire une page scannée, mais ils produisent aussi des informations que le document ne contient pas. Dans un contexte médical, ces erreurs, appelées hallucinations, ne sont pas acceptables sans contrôle.

Pour y répondre, ce mémoire propose quatre contributions : un guideline d'annotation validé par un médecin, un Gold Standard construit à la main sur 49 dossiers de la cohorte APROCCHSS, un pipeline d'extraction multimodale guidé par ce guideline, et TRACE, une chaîne de vérification qui détecte et réduit les hallucinations après l'extraction. Une partie de ces travaux a été présentée lors de trois communications : un article à la conférence JDSE~2026, une présentation orale à PFIA~2026 et un poster à la Journée des Doctorants \& Post-doctorants de l'IHU SEPSIS et de l'OI MICROBES.

Ce projet m'a permis de travailler à l'interface entre l'intelligence artificielle et la médecine, au contact de cliniciens et de chercheurs. Il m'a surtout appris qu'en santé, la fiabilité d'un système compte autant que ses performances : un résultat n'a de valeur que si l'on peut dire d'où il vient et à quel point on peut s'y fier. C'est ce principe qui a guidé chacun des choix présentés dans ce mémoire.

Le lecteur trouvera d'abord un état de l'art sur l'extraction d'information clinique et les hallucinations des modèles de langage, puis la description du guideline et du Gold Standard, celle du pipeline d'extraction, celle de la chaîne TRACE, et enfin l'ensemble des expérimentations et des résultats.
